\begin{textsample}{NEG Dim 7 – human – Score -2.00 – mg/mg\_ef\_linguagens\_ed\_fis\_8\_p3.txt}  \label{ex:f7_neg_018}
Nos anos finais do ensino fundamental já é possível investir em estratégias e metodologias que estimulem capacidades reflexivas, críticas e criativas, essenciais para a formação cidadã.
As estruturas cognitivas se fortalecem e se consolidam podendo ser veículo para a constituição de conceitos, princípios e valores.
Nessa etapa é desejável o investimento em metodologias que valorizem os conhecimentos técnicos e táticos vinculados às diversas práticas corporais, assim como estratégias que possibilitem seu treinamento e refinamento.
As conquistas motoras e sensoriais devem estar a serviço de ideais de nobreza e retidão, \textbf{colocando} a corporeidade disponível para ações que gerem benefícios coletivos e que promovam valores inclusivos e recursos para o exercício do protagonismo social.

% matched lemmas: colocar
\end{textsample}
